\definecolor{rustcolor}{HTML}{DEA584}
\definecolor{corecolor}{HTML}{4A90D9}
\definecolor{storecolor}{HTML}{50C878}
\definecolor{netcolor}{HTML}{9B59B6}
\definecolor{clicolor}{HTML}{E74C3C}
\definecolor{opcolor}{HTML}{F39C12}
\definecolor{vccolor}{HTML}{1ABC9C}
\definecolor{bglight}{HTML}{F8F9FA}

\begin{tikzpicture}[
    node distance=1.5cm and 1cm,
    opnode/.style={rectangle split, rectangle split parts=4,
                   draw, fill=white, rounded corners=3pt, thick,
                   font=\footnotesize\ttfamily, align=left,
                   minimum width=3.5cm},
    orderidx/.style={circle, inner sep=1.5pt, font=\small\sffamily\bfseries},
    parentlink/.style={-{Stealth[length=3mm, width=2.5mm]}, very thick, draw=opcolor, dashed}
]

% Title
\node[font=\Large\sffamily\bfseries] (title) at (4, 11) {equack-core: Operation Structure \& Causal DAG};

% Legend
\node[draw=gray, fill=white, rounded corners, align=left, font=\scriptsize\sffamily, thick] at (13.5, 10.5) {
    \textbf{Topological Order Tie-breaking:}\\
    1. Parents strictly before children\\
    2. Lowest HLC (physical\_ms, logical)\\
    3. Lowest Op-ID (BLAKE3)
};

% ============================================================================
% MAIN DIAGRAM STRUCTURE
% ============================================================================
% The user provided a snippet that seems to replace the entire tikzpicture content
% but also included a new \begin{tikzpicture}. I'm assuming the intent is to
% replace the content within the existing tikzpicture environment.
% Also, 'policycolor' is used in the new code but not defined. Assuming it's
% meant to be 'vccolor' or a new color definition is missing.
% For now, I'll define policycolor as vccolor.
\definecolor{policycolor}{HTML}{1ABC9C} % Assuming policycolor is equivalent to vccolor

\tikzset{
    opnode/.style={rectangle split, rectangle split parts=4,
                   draw, fill=white, rounded corners=3pt, thick,
                   font=\footnotesize\ttfamily, align=left,
                   minimum width=3.5cm},
    orderidx/.style={circle, inner sep=1.5pt, font=\small\sffamily\bfseries},
    parentlink/.style={-{Stealth[length=3mm, width=2.5mm]}, very thick, draw=opcolor, dashed}
}

    % ========================================================================
    % NODE DEFINITIONS
    % Nodes placed with much broader absolute horizontal and relative vertical
    % spacing to cleanly avoid multi-part text node overlaps.
    % ========================================================================

    \node[opnode, draw=corecolor!80!black] (op1) at (0, 8) {
        % Base / Header
        \textbf{Op-ID: A1B2...} \nodepart{two}
        \texttt{Header:} \\
        \texttt{- parents: [ ]} \\
        \texttt{- hlc:  (1000, 0, N1)} \\
        \texttt{- author: PK\_A} 
        % Payload
        \nodepart{three}
        \texttt{Payload: Trust(IssuerKey)}
        % Signature
        \nodepart{four}
         \texttt{Sig: Ed25519(...)}
    };
    
    % Event order circle
    \node[orderidx, fill=black, text=white, right=3pt] at (op1.north west) {1};

    \node[opnode, draw=policycolor!80!black] (op2) at (4, 4) {
        \textbf{Op-ID: C3D4...} \nodepart{two}
        \texttt{Header:} \\
        \texttt{- parents: [ A1B2... ]} \\
        \texttt{- hlc:  (1005, 0, N1)} \\
        \texttt{- author: PK\_A} 
        \nodepart{three}
        \texttt{Payload: Credential(VC)}
        \nodepart{four}
        \texttt{Sig: Ed25519(...)}
    };
    \node[orderidx, fill=black, text=white, right=3pt] at (op2.north west) {2};

    \node[opnode, draw=storecolor!80!black] (op3) at (8, 8) {
        \textbf{Op-ID: 7G8H...} \nodepart{two}
        \texttt{Header:} \\
        \texttt{- parents: [ A1B2... ]} \\
        \texttt{- hlc:  (1008, 0, N3)} \\
        \texttt{- author: PK\_C} 
        \nodepart{three}
        \texttt{Payload: Data(MVReg=Val1)}
        \nodepart{four}
        \texttt{Sig: Ed25519(...)}
    };
    \node[orderidx, fill=black, text=white, right=3pt] at (op3.north west) {3};

    \node[opnode, draw=policycolor!80!black] (op4) at (0, 0) {
        \textbf{Op-ID: E5F6...} \nodepart{two}
        \texttt{Header:} \\
        \texttt{- parents: [ C3D4... ]} \\
        \texttt{- hlc:  (1010, 0, N2)} \\
        \texttt{- author: PK\_B} 
        \nodepart{three}
        \texttt{Payload: Grant(Key, Scope)}
        \nodepart{four}
        \texttt{Sig: Ed25519(...)}
    };
    \node[orderidx, fill=black, text=white, right=3pt] at (op4.north west) {4};

    \node[opnode, draw=storecolor!80!black] (op5) at (6, 0) {
        \textbf{Op-ID: 9I0J...} \nodepart{two}
        \texttt{Header:} \\
        \texttt{- parents: [ E5F6..., 7G8H... ]} \\
        \texttt{- hlc:  (1015, 0, N2)} \\
        \texttt{- author: PK\_B} 
        \nodepart{three}
        \texttt{Payload: Data(MVReg=Val2)}
        \nodepart{four}
        \texttt{Sig: Ed25519(...)}
    };
    \node[orderidx, fill=black, text=white, right=3pt] at (op5.north west) {5};

    % ========================================================================
    % DAG EDGES (Parent references)
    % Routed with explicit angles to prevent them slicing through the text nodes
    % ========================================================================
    \begin{scope}[on background layer]
        % op2 -> op1
        \draw[parentlink] (op2.north west) to[out=135, in=315] (op1.south east);
        \node[font=\tiny\sffamily, text=black!50] at (2, 5.8) {references};

        % op3 -> op1
        \draw[parentlink] (op3.west) to[out=180, in=0] (op1.east);

        % op4 -> op2
        \draw[parentlink] (op4.north east) to[out=45, in=225] (op2.south west);

        % op5 -> op4
        \draw[parentlink] (op5.west) to[out=180, in=0] (op4.east);
        
        % op5 -> op3
        \draw[parentlink] (op5.north east) to[out=45, in=225] (op3.south);
    \end{scope}

% Epoch Brackets (Right side highlighting Replay Engine logic)
\draw[thick, draw=corecolor] (12, 7) -- (12, 3) node[midway, right=10pt, align=left, font=\footnotesize\sffamily] {
    \textbf{Auth Epoch Builder}\\
    Evaluates Trust, Credentials,\\
    and Grants in order to gate\\
    upcoming Data writes.
};

\draw[thick, draw=storecolor] (12, 2.5) -- (12, -1) node[midway, right=10pt, align=left, font=\footnotesize\sffamily] {
    \textbf{Deny-Wins Gate}\\
    Data Ops accepted only if\\
    VC valid, unrevoked, inside\\
    time window \mbox{[nbf, exp)} and\\
    intersecting scope.
};

\end{tikzpicture}
